Une planète orbite à \qty{\v{dp}}{UA} du Soleil et possède un albédo de \num{\v{al}}. À la distance de la Terre (\qty{1}{UA}), l'intensité du rayonnement solaire est de \qty{1361}{W/m^2}.

\begin{enumerate}[label=\alph*)]
    \item Déterminez l'intensité du rayonnement solaire à la distance de la planète.
    \item Déterminez la température d'équilibre de la planète en l'assimilant à un corps noir sans effet de serre.
    \item Si on observe à sa surface une température moyenne beaucoup plus élevée que ce résultat, que peut-on conclure sur son atmosphère ?
\end{enumerate}
